\section{Related Work}\label{sec:related}
\noindent\textbf{EEG foundation models.} Early self-supervised EEG models used contrastive or
autoregressive objectives (BENDR \cite{kostas2021bendr}, Neuro-GPT \cite{cui2024neurogpt}).
Most recent models pretrain a transformer on TUEG with a masked-prediction objective: CBraMod
\cite{wang2025cbramod}, CSBrain \cite{csbrain} and REVE \cite{reve2025} reconstruct the raw
waveform of masked patches; LaBraM \cite{jiang2024labram} predicts the codes of a
vector-quantized tokenizer; EEGPT \cite{wang2024eegpt} combines reconstruction with latent
alignment; NeuroLM \cite{jiang2025neurolm} maps EEG to language tokens; BIOT
\cite{yang2023biot} handles heterogeneous biosignals. These models differ mainly in
architecture, scale and data, while their targets are inherited from vision and language. A
waveform target asks the model to predict sample-level details, including phase, in addition
to spectral structure; some of this variation may remain uncertain given the context. A
learned codebook target is defined by a tokenizer rather than by a physiological quantity. Our
target directly measures band power, a spectral quantity relevant to many EEG labels.

\noindent\textbf{Spectral information in EEG pretraining.} Spectral features enter several
models as inputs or tokenizers: BIOT tokenizes STFT magnitudes, CBraMod adds a
frequency-domain patch embedding, LaBraM's tokenizer reconstructs the Fourier amplitude and
phase, BioSerenity-E1 \cite{bettinardi2025bioserenity} and BandVQ \cite{sukhbaatar2026bandvq}
quantize log spectra, and BrainBERT \cite{wang2023brainbert} reconstructs intracranial
spectrograms. Concurrently, FAME \cite{yu2026fame} decodes band-averaged log power of the whole
sequence to counter the low-frequency bias of masked autoencoders, and MANAS-2
\cite{kulkarni2026manas2} adds band log-power to waveform reconstruction. In these works
spectral information is an input, a quantized code or an auxiliary term, and its role as the
target is neither motivated nor isolated. We use band power of the masked patches as the sole
target, derive it from a likelihood argument, and isolate its effect with matched controls.

\noindent\textbf{Feature and latent prediction.} In vision and speech, predicting fixed
hand-crafted features of masked inputs (HOG in MaskFeat \cite{wei2022maskfeat}, MFCC clusters in
HuBERT \cite{hsu2021hubert}) can outperform predicting the raw signal, and latent-predictive
methods (I-JEPA \cite{assran2023ijepa}, data2vec \cite{baevski2022data2vec}, EEG2Rep
\cite{foumani2024eeg2rep}) predict learned representations but need anti-collapse mechanisms
such as SIGReg \cite{balestriero2025lejepa}. Band power plays the role of such a feature for
EEG: it is a physiological quantity relevant to many labels, and we use the fixed target
without an explicit anti-collapse regularizer.
